% =====================================================================
\section{Lighting, Materials and Surface Detail}\label{sec:shading}
\pipestrip{6}
\subsection{Lambert and Blinn--Phong: diffuse and shininess}
\autofig{diagrams}{B3_lambert}
\autofig{figures_lab}{materials_blinnphong_row}
For a light with direction $\vect L$ (towards the light), normal $\vect N$, view vector $\vect V$ and halfway vector $\vect H=\mathrm{normalize}(\vect L+\vect V)$, the classic model is
\begin{equation}
I=k_d\,\max(\vect N\!\cdot\!\vect L,0)\;+\;k_s\,\max(\vect N\!\cdot\!\vect H,0)^{n}.
\end{equation}
The first term is \emph{Lambert's law}: received energy follows $\cos\theta=\vect N\!\cdot\!\vect L$, scattered evenly. The second is the highlight; the \emph{shininess} $n$ sets its width, which falls to half at $\alpha_{1/2}=\arccos 2^{-1/n}$: $9.4^\circ$ for $n=16$, $4.7^\circ$ for $n=64$, $2.4^\circ$ for $n=256$. This is the \code{uPbr==0} branch of Listing~\ref{lst:Brdf} (key \texttt{B}); \code{uSpecular} is $k_s$, \code{uShininess} is $n$.

\subsection{Cook--Torrance with GGX: roughness and metalness}
\autofig{diagrams}{C3_specular_grid}
The default model treats a surface as microscopic mirrors (\emph{microfacets}); with roughness $\rho$, $\alpha=\rho^2$ and
\begin{equation}
f_s=\frac{D(\vect H)\,G(\vect N,\vect V,\vect L)\,F(\vect V,\vect H)}{4(\vect N\!\cdot\!\vect V)(\vect N\!\cdot\!\vect L)},
\end{equation}
\begin{equation}
D=\frac{\alpha^2}{\pi\bigl((\vect N\!\cdot\!\vect H)^2(\alpha^2-1)+1\bigr)^2},\qquad
G=\prod_{x\in\{V,L\}}\frac{\vect N\!\cdot\!x}{(\vect N\!\cdot\!x)(1-k)+k},\ \ k=\frac{(\rho+1)^2}{8},\qquad
F=F_0+(1-F_0)(1-\vect V\!\cdot\!\vect H)^5 .
\end{equation}
$D$ is the GGX normal distribution (facets facing $\vect H$, with the long tail that makes highlights natural), $G$ the Smith--Schlick masking term (facets hide each other at grazing angles) and $F$ the Schlick Fresnel term (every surface mirrors at grazing angles). The base reflectivity is $F_0=\mathrm{mix}(0.04,\text{albedo},m)$ with metalness $m$: dielectrics reflect 4\,\% white, metals their own colour; the diffuse part is $(1-F)(1-m)$ times Lambert. Listing~\ref{lst:Brdf} is the whole function.
\lstBrdf
\paragraph{Shininess $\leftrightarrow$ roughness.} Materials that give only a shininess get $\rho=\mathrm{clamp}(1-\log_2(n)/8,\,0.08,\,0.95)$ (\code{pbrRoughness}, \code{Material.h}); others set roughness explicitly. Figure~\ref{fig:matgrid} shows roughness against metalness under one light.
\begin{figure}[h]
\centering
\begin{subfigure}{0.58\linewidth}\piclab{material_grid}{\linewidth}\caption{roughness (columns) $\times$ metalness (rows)}\end{subfigure}
\begin{subfigure}{0.40\linewidth}\picdia{ggx_curves}{\linewidth}\caption{$D$ for several $\rho$; Blinn--Phong lobes}\end{subfigure}
\caption{Material response. The highlight narrows and brightens as roughness falls; metals tint the highlight, dielectrics do not.}\label{fig:matgrid}
\end{figure}

\subsection{Light sources and attenuation}
\begin{itemize}
  \item \textbf{L0, moon:} directional, colour $(0.55,0.62,0.80)$, intensity 0.95.
  \item \textbf{L1, beacon:} a point light at the bunker's beacon, colour $(1,0.25,0.10)$, attenuation $1/(1+0.35\,d+0.45\,d^2)$, intensity pulsing $0.6$--$2.0$ at 1\,Hz.
  \item \textbf{L2, UAV searchlight:} a spot with smooth cone between $9^\circ$ and $14^\circ$: $c=\mathrm{clamp}\frac{\cos\gamma-\cos14^\circ}{\cos9^\circ-\cos14^\circ}$, radiance $\propto c^2/(1+0.05\,d+0.03\,d^2)$, intensity 5.
  \item \textbf{Practical lights:} any emissive part declared with \code{asLight(material, intensity, range, cone)} is also a light. The renderer ranks them by distance over range and uploads the \textbf{16} best per view (\code{gatherPracticals}). Falloff is a windowed inverse square (Listing~\ref{lst:Practical}): $f=\mathrm{saturate}(1-(d/R)^4)^2/(1+d^2)$, exactly 0 at the range $R$ (no popping), optionally cut by a spot cone.
  \item \textbf{Dynamic explosion lights:} flashes, burning wrecks and muzzle flashes join the same list and rank ahead of lamps (score halved).
\end{itemize}
\lstPractical
\begin{figure}[h]
\centering
\picdia{attenuation}{0.7\linewidth}
\caption{Attenuation curves of the beacon, the searchlight and a practical lamp (computed from the constants above).}
\end{figure}

\begin{figure}[h]
\centering
\picsheet{sheet_lighting}{\linewidth}
\caption{\sheetcaption{sheet_lighting}}
\end{figure}
\subsection{Shadows: cascades and PCF}
\sheetfig{s07_attenuation}
\sheetfig{s07_charts}
The moon casts shadows with two orthographic cascades in a depth-texture array. The \emph{near cascade} covers the view-frustum slice from the near plane to a split distance (10--14\,m by preset) and follows the camera; the \emph{world cascade} covers a 39\,m sphere around the base, never moves, and is snapped to whole texels so edges do not shimmer (T34). Its static casters are rendered once into a cache; only \code{dynamic} nodes (vehicles, UAVs, debris) are redrawn each frame. The cache key is the explode value plus the foliage setting. Listing~\ref{lst:MoonShadow} blends the cascades over the last 20\,\% of the split distance. Each lookup uses a normal-offset bias and a rotated Poisson disc of 8 (HIGH) or 16 (ULTRA/EPIC) hardware-PCF taps; the fraction of lit taps gives a soft penumbra. The searchlight has its own perspective shadow map along its axis.
\lstMoonShadow
\begin{figure}[h]
\centering
\piclab{shadow_cascades}{0.8\linewidth}
\caption{Moon shadow cascades (colour = cascade).}
\end{figure}

\subsection{Ambient light, indoor mask, thermal heat}
Ambient is a sky/ground hemisphere. The bunker must receive no sky or moonlight: \code{WalkSpace} supplies up to eight indoor boxes and Listing~\ref{lst:Indoor} computes a smooth mask. Inside it ambient is multiplied by $0.012$, moon, beacon and searchlight by $1-\text{mask}$, fog by $1-0.85\,\text{mask}$, and practical lights carry an \emph{inside} flag so a tunnel lamp does not light the surface above.
\lstIndoor
The thermal view reuses the shader with apparent temperature $h=h_0+0.04\min(E,1.5)$ (intrinsic heat plus a little capped absorbed light), minus $0.05$ times the wetness on passive surfaces.

\subsection{Procedural albedo patterns}
There are no image textures. The \code{Surface} enum selects a pattern evaluated from world position with value, fractal ($\mathrm{fbm}$) or cellular noise: concrete, paint chips, bark, lichen, leaves, and the set of Table~\ref{tab:surfaces}. Patterns are triplanar: coordinates are $(x,z)$ on up-facing surfaces and $(z,y)$ or $(x,y)$ on side faces by dominant normal axis, so a pattern continues across a scaled cube without UV unwrap.
\begin{table}[h]
\centering\small
\caption{Procedural surfaces (\code{Surface} values 12--18).}\label{tab:surfaces}
\begin{tabular}{@{}rllp{8.4cm}@{}}
\toprule Id & Name & Used on & Pattern\\\midrule
12 & Corrugated & containers, huts, radar shelter & ribs every 7.6\,cm along world $Y$, dirt runs, rust at the bottom\\
13 & Hesco & HESCO bastion walls & welded wire grid (7.6\,cm) over sand-bag bulges\\
14 & Hazard & bunker door portal, barriers, signs & yellow/black $45^\circ$ stripes with wear\\
15 & Screen & bunker consoles, control panel & dark phosphor, grid, rotating sweep, blips, text bars (emissive)\\
16 & Asphalt & the gate road & aggregate speckle, cracks, polished wheel lanes\\
17 & Camo & tank, helicopter, AA gun, nets & three-blob woodland camouflage\\
18 & SteelPlate & blast doors, armour, radar door & 0.6\,m panels, seam lines, rivet rows, worn edges\\\bottomrule
\end{tabular}
\end{table}
\begin{figure}[h]
\centering
\piclab{surface_patterns}{\linewidth}
\caption{The procedural surfaces on flat boxes (in-program render): concrete, corrugated, HESCO, hazard, screen, asphalt, camo, steel plate.}
\end{figure}

\subsection{Surface relief: procedural bump}
\code{surfaceHeight(q)} in \code{ShaderSources.cpp} gives each procedural material a height field in the same triplanar coordinates as its albedo (concrete pores 1.5\,mm, bark ridges, rock, planks, corrugated ribs 9\,mm, HESCO mesh and bags, asphalt
aggregate, camo paint chips, steel seams $-3$\,mm and rivets $+2$\,mm). Listing~\ref{lst:Bump} shows \code{surfaceBump(P,N)}, which takes central differences $\partial h/\partial u,\partial h/\partial v$ ($\epsilon=4$\,mm) and tilts the normal,
$\vect N_b=\mathrm{normalize}\bigl(\vect N-(h_u\vect T+h_v\vect B)\,k\bigr)$, where $\vect T,\vect B$ is the tangent frame of the dominant face and $k$ fades out from 12 to 40\,m (no far shimmer). Because the Gouraud and Flat paths evaluate lighting before the fragment
stage, the light terms are \emph{re-aimed} by the Lambert ratio
\begin{equation}
k_L=\mathrm{clamp}\!\left(\frac{\max(\vect N_b\!\cdot\!\vect L,0)+0.03}{\max(\vect N\!\cdot\!\vect L,0)+0.03},\,0.3,\,2\right),
\end{equation}
applied to the moon and searchlight diffuse terms ($k_L$) and specular terms ($k_L^2$); a cavity factor $\mathrm{clamp}(0.6+0.4\,\vect N_b\!\cdot\!\vect N+0.5(N_{b,y}-N_y),0.6,1.15)$ darkens pits that face away from the sky, and the bumped normal also feeds the SSR surface buffer. The bump is active
only in HDR on HIGH and above (\code{uGroundDetail}); screens, foliage and the ground have their own code. \lstBump
\begin{figure}[h]
\centering
\piclab{bump_on_off}{\linewidth}
\caption{Bump off (left) and on (right) for corrugated metal, HESCO and steel plate: ribs, mesh and rivets catch the light (in-program render).}
\end{figure}
The planned \emph{reflection probe} (a low-resolution sky cubemap for metal and glass) is \textbf{not} implemented; reflections still come from SSR and the analytic sky.

